\section{Checks for the annular cap comparison}
\label{app:annular-cap-comparison}

The algebraic constants in \cref{thm:annular-cap-comparison}
satisfy
\[
 \frac{(53/40)^2+8(53/40)}9=\frac{19769}{14400}<\frac{11}{8},
 \qquad
 \frac{32\cdot27\cdot6^3}{105\cdot4^7(314/100)}<\frac7{200},
 \qquad \frac{5}{3}\frac7{200}=\frac7{120}.
\]
Here \(4^7=16^{7/2}\). These comparisons need only integer
arithmetic after cross multiplication.

The following example checks the effect of the two volume
bounds on the same shell. Let \(Y=2hU\), where \(U\) is the
set of twenty-four unit \(D_4\) roots and \(h\) is either
\(1\) or \(10031/10000\). These are admissible shells since
each pair of distinct unit roots has distance at least \(1\)
and \(2h<\sqrt6\). The cap plane distance is \(h\).
Each root has eight other roots at inner product \(1/2\),
so there are ninety-six such unordered pairs. All other
pairs have inner product at most \(0\), and cannot cut the
same point of \(B_\rho\), since its norm would be at least
\(\sqrt2h>\rho\).

Put \(\alpha=\arccos(h/r)\) for \(r\in\{r_0,\rho\}\).
Integration of perpendicular three-balls gives the single
cap volume
\[
 S_r(2h)=\frac{4\pi}{3}
 \left(\frac{3r^4}{8}\alpha-
       \frac{h(5r^2-2h^2)}8\sqrt{r^2-h^2}\right).
\]
For completeness, use polar coordinates in the plane of
the two cap normals, at angle \(\pi/3\) to one another.
The angle of the projection runs from \(\pi/3-\alpha\)
to \(\alpha\). On the half from \(\pi/6\) to \(\alpha\),
the first normal requires projected radius at least
\(h/\cos\phi\), and the other half is its reflection.
The two perpendicular coordinates contribute a disk of
area \(\pi(r^2-s^2)\) at projected radius \(s\).
Integrating \(s\) from \(h/\cos\phi\) to \(r\), and doubling,
gives the pair volume
\[
 \Pi_r=\frac{\pi}{2}\bigl(G_r(h,\alpha)-G_r(h,\pi/6)\bigr),
 \qquad
 G_r(h,\phi)=r^4\phi-2r^2h^2\tan\phi+
                    h^4\left(\tan\phi+\frac{\tan^3\phi}{3}\right).
\]
Indeed \(G_r'(h,\phi)=(r^2-h^2\sec^2\phi)^2\).
Thus \(C_r=24S_r(2h)\), \(P_r=96\Pi_r\), and the value of
\(T(Y)\) uses the exact two-term inclusion--exclusion formula
at \(r_0\).

\begin{table}[ht]
\centering
\caption{Rational enclosures for the truncated volume and
the expressions \(T+(77/100)L\) from the two annular bounds.
The positive term \(g(\rho)\) is omitted from both expressions.}
\label{tab:annular-cap-check}
\begin{tabular}{@{}lccc@{}}
\toprule
\(h\) & \(T(Y)\) & \cref{eq:annular-cap-one-pair} & \cref{eq:annular-cap-exact-pairs}\\
\midrule
\(1\) & \( (7.9069,\,7.9070) \) & \( (7.9146,\,7.9147) \) & \( (7.9360,\,7.9361) \)\\
\(10031/10000\) & \( (7.9962,\,7.9963) \) & \( (8.0008,\,8.0009) \) & \( (8.0291,\,8.0293) \)\\
\bottomrule
\end{tabular}
\end{table}

The script \texttt{anc/cap\_comparison.py} encloses square
roots on a grid of denominator \(10^{24}\) by integer square
root. It evaluates \(\arctan(\sqrt{r^2-h^2}/h)\) on the
endpoint fractions with the first forty-eight terms of the
alternating series and the first omitted term as an error
bound. It uses the same series for the Machin formula in
\cref{app:annular-constants}. All other operations are on
rational intervals.\footnote{The argument of each arctangent
is in \([0,1)\) here, so the error bound follows directly
from alternating terms of decreasing magnitude.}
For \(h=10031/10000\), even the bound using pair data only
at \(r_0\) proves volume greater than \(8\) for every
exterior extension, although \(T(Y)<8-1/28600\).
The \(h=1\) row records the loss of this criterion at the
regular cell, whose volume is exactly \(8\).

\section{A smaller certified level at twenty-six neighbours}
\label{app:profiles-26}

Use the seven bins with lower endpoints
\[
 (2,\ 2.0161,\ 2.05,\ 2.1,\ 2.15,\ 2.2,\ 2.444)
\]
from the \(M=26\) part of \cref{prop:levels}, with types
\(\mathsf{A,A,B,B,B,B,F}\). Let \(n=(n_0,\ldots,n_6)\)
be the vector of occupancies. If \(T(Y)\le8\), then
\(n_0\le24\), the total is \(26\), and the cumulative
occupancies through bins \(1,\ldots,5\) are at least
\((15,20,22,23,24)\), by the radial and cardinality
certificates in that proof. In particular \(n_6\le2\).

The previous certified kernel bounds the union of caps \(U(Y)\)
above by
\[
 F(n)=\frac{t}{2}+\sum_{b=0}^6 n_bq_b
   +\sum_{\mathsf{st}}N_{\mathsf{st}}c_{\mathsf{st}}
   +\sum_{\mathsf{stu}}N_{\mathsf{stu}}c_{\mathsf{stu}},
\]
with \(t=8.282667610\). The point coefficients, in bin order, are
\[
 \begin{array}{rrrr}
  0.10730785 & 0.10730646 & -0.05347062 & -0.05363574\\
  -0.05425642 & -0.05347062 & -0.03946852 &
 \end{array}
\]
Each \(N_w\) counts unordered subsets with type word \(w\).
If the type counts are \(n_{\mathsf A},n_{\mathsf B},n_{\mathsf F}\),
then \(N_w=\prod_{\mathsf s}\binom{n_{\mathsf s}}{m_{\mathsf s}(w)}\),
where \(m_{\mathsf s}(w)\) is the multiplicity of \(\mathsf s\) in \(w\).
The constants \(c_w\) in this computation are the entries below
plus \(10^{-7}\) in each case.
\[
\begin{array}{c|rrr}
w & \mathsf{AA} & \mathsf{AB} & \mathsf{AF}\\
\hline
c_w-10^{-7} & -0.02327176 & -0.01036330 & -0.01503736\\
w & \mathsf{BB} & \mathsf{BF} & \mathsf{FF}\\
\hline
c_w-10^{-7} & -0.001713338 & -0.004512281 & -0.003681549
\end{array}
\]
\[
\begin{array}{c|rrr}
w & \mathsf{AAA} & \mathsf{AAB} & \mathsf{AAF}\\
\hline
c_w-10^{-7} & .001497925 & .0009598990 & .001322105\\
w & \mathsf{ABB} & \mathsf{ABF} & \mathsf{AFF}\\
\hline
c_w-10^{-7} & .0005966748 & .001043879 & .001145076\\
w & \mathsf{BBB} & \mathsf{BBF} & \mathsf{BFF}\\
\hline
c_w-10^{-7} & .0004094797 & .0006564131 & .001043982
\end{array}
\]
The type \(\mathsf{FFF}\) has multiplicity zero.
The constants are rounded upward from the archived checks
of the individual kernel inequalities.\footnote{The input is
\href{https://github.com/creelie/d4-voronoi-cells/blob/v3.4.0/multi_cap/runs/combo26_level_check.log}{the v3.4.0
certificate log} and the array \texttt{m} in
\texttt{multi\_cap/radial\_certificates/combo26\_level\_d8.npz}.
The \(q_b\) combine the recorded point term for the type with
the upward-rounded bin bound \(m_b+2\cdot10^{-6}\) of the
checker. The rational postprocessing here takes the
kernel inequalities as inputs. For related methods of
certifying semidefinite packing bounds, see \cite{DLM21}.}

There is a second bound, independent of the kernel.
If \(d_b\) is the lower endpoint of bin \(b\), monotonicity
of the single cap gives
\[
 U(Y)\le \sum_b n_b S_{r_0}(d_b).
\]
Taking the minimum of this expression and \(F(n)\) before
maximising over \(n\) gives a smaller level. The script
\texttt{anc/cap\_profiles\_26.py} enumerates the \(17\,845\)
weak compositions satisfying the cumulative constraints.
It rounds the cap bounds upward using the interval procedure
in \cref{app:annular-cap-comparison} and finds that
\[
 \max_n\min\left\{F(n),\ \sum_b n_b S_{r_0}(d_b)\right\}
 <3.288 .
\]
The upper bounds in the script are largest at
\((20,0,5,0,0,0,1)\). They give \(T(Y)>8\) on \(12\,378\)
of the profiles. The profile \((15,0,9,0,0,0,2)\) that
previously gave the largest kernel value has a cap sum less
than \(3.027\) and so already gives \(T(Y)>8\).
If \(T(Y)>8\), then \(U(Y)<9\pi^2/8-8<3.288\) directly.
Thus \(U(Y)<3.288\) holds at this neighbour count without
any assumption on \(T(Y)\).

The largest level in \cref{tab:levels} is consequently
at \(M=25\). In that case the two-bin certificate has
upward bounds \(t=1.819065983\), \(m_0=0.097085101\)
and \(m_1=0.087814019\), with zero multiplier
for the radial side condition.\footnote{These are
rounded upward from the coefficient array in
\texttt{multi\_cap/radial\_certificates/density\_25.json}
and the value of \(t\) in \texttt{gap\_closure/density/check\_25.log}
in the same archive.}
Since \(m_0>m_1\), at most twenty-three neighbours
within \(2.0161\) gives
\[
 U(Y)\le t/2+23m_0+2m_1<3.3182.
\]
Thus only twenty-five-neighbour shells with exactly
twenty-four within \(2.0161\) can exceed the level \(3.3182\)
after these refinements. The sharp level \(9\pi^2/8-8\)
would require a further improvement.
